% =============================================================================
% ADDENDUM: Formal Welfare Derivation for "Discrimination from Soft Signals"
% =============================================================================
% This is a standalone reference document containing the full formal welfare
% derivation. It is NOT part of the submission to Theory and Decision. The main
% paper carries welfare only as a future-work note in the Conclusion; this
% addendum develops that direction in full — the welfare object, the telescoping
% decomposition, the within-group coarseness floor, the three-regime framing,
% and the worked policy comparison — for reference.
%
% Companion file: manuscript.tex
% Verification: verify_discrimination.py CHECK_86 through CHECK_90 and
%               ANALYTIC_10, ANALYTIC_11
% =============================================================================

\documentclass[12pt,a4paper]{article}

%% Packages
\usepackage[a4paper, margin=1.1in]{geometry}
\usepackage{amsmath, amssymb, amsthm}
\usepackage{setspace}
\usepackage{booktabs}
\usepackage{array}
\usepackage{hyperref}
\usepackage{natbib}
\usepackage{graphicx}
\usepackage{enumitem}
\usepackage{titlesec}
\usepackage{parskip}
\usepackage[expansion=false]{microtype}
\usepackage{xcolor}

%% Spacing and layout
\onehalfspacing
\setlength{\parindent}{0pt}
\setlength{\parskip}{8pt}

%% Theorem environments
\newtheorem{proposition}{Proposition}
\newtheorem{theorem}[proposition]{Theorem}
\newtheorem{lemma}[proposition]{Lemma}
\newtheorem{corollary}[proposition]{Corollary}
\newtheorem{definition}[proposition]{Definition}
\newtheorem{remark}[proposition]{Remark}
\theoremstyle{definition}
\newtheorem{example}[proposition]{Example}
\numberwithin{proposition}{section}

%% Title formatting
\titleformat{\section}{\large\bfseries}{\thesection.}{0.5em}{}
\titleformat{\subsection}{\normalsize\bfseries}{\thesubsection.}{0.5em}{}
\titleformat{\subsubsection}{\normalsize\itshape}{\thesubsubsection.}{0.5em}{}

%% Hyperref
\hypersetup{colorlinks=true, linkcolor=blue!60!black,
            citecolor=blue!60!black, urlcolor=blue!60!black}

%% Custom commands
\newcommand{\E}{\mathbb{E}}
\newcommand{\R}{\mathbb{R}}
\newcommand{\G}{\mathcal{G}}
\newcommand{\M}{\mathcal{M}}
\newcommand{\C}{\mathcal{C}}
\newcommand{\A}{\mathcal{A}}

\begin{document}

\begin{center}
{\Large\bfseries Addendum: Formal Welfare Derivation}\\[0.4em]
{\normalsize Companion to ``Stereotype Rigidity in Social Inferences:
A `Mimic Men' Dilemma''}\\[0.5em]
{\normalsize Reference document --- not for submission}\\[1em]
\end{center}

\section*{Scope of this document}

The main paper carries the welfare analysis only as a brief future-work
note in the Conclusion: the natural benchmark is a first-best in which
social engagement conditions on an individual's true class rather than
on a group marker, and relative to it the loss the trap generates
separates into a component driven by the group's standing and a
component that survives even when that standing is accurate. This
document develops that direction in full: the welfare object $W$
defined as a discounted population integral; the telescoping
construction by which the four components arise; the integral
expressions for each component, including the within-group coarseness
floor that is the load-bearing object here; the three-regime case
framing; the proof of additivity; and the worked policy comparison
under each regime. The material is verified symbolically in
\texttt{verify\_discrimination.py} as CHECK\_86 through CHECK\_90
and ANALYTIC\_10, ANALYTIC\_11. It is a reference document, organised
around a single headline result (welfare dominance of mobility above
$\gamma^*$) and one companion negative result (irreducibility of the
within-group informational loss); it is not part of the submission.

\section{Welfare object and decomposition}
\label{sec:welfare-object}

We define social welfare $W$ as the population-aggregate expected
utility from social engagement, summed across encounters within a
generation and discounted across generations:
\begin{equation}
  W \;=\; \E\!\left[\,
  \sum_{t \geq 0} \beta^t
  \int\!\!\!\int U(c_i, c_j) \cdot \mathbf{1}\!\left[\text{engage}\right]
  \, dF^t(c_i)\, dF^t(c_j)
  \,\right]
  \label{eq:welfare-W}
\end{equation}
where $\beta \in (0,1)$ is the intergenerational discount factor,
$F^t$ is the equilibrium class distribution in generation $t$, and
$U(c_i, c_j)$ is the bilateral utility from the main paper's
equation~(15). Welfare is utilitarian: a realised encounter
contributes $U(c_i, c_j)$ to surplus, while a refused encounter
contributes the outside option $\bar{u}$, normalised to zero.

The first-best welfare $W^*$ is the value of $W$ achieved under three
simultaneous conditions: (i) every mutually beneficial encounter
takes place---$U(c_i, c_j) > \bar{u}$ implies engagement; (ii) social
engagement conditions on each individual's \emph{true} class $c_i$
rather than on the group marker $m_i$, so that no surplus is foregone
to the coarseness of group-level information; and (iii) no shadow
information market is active. Condition~(ii) is the substantive
benchmark: it does not ask that beliefs about the group be accurate,
but that engagement need not run through the group at all. The total
welfare loss is the gap between this first-best and equilibrium
welfare,
\begin{equation}
  \mathcal{L} \;=\; W^* - W^{\mathrm{eq}}.
\end{equation}

To decompose $\mathcal{L}$ we construct a sequence of counterfactual
welfares that interpolate between $W^*$ and $W^{\mathrm{eq}}$, each
relaxing one of the conditions above:
\begin{align*}
  W_0 &\equiv W^* && \text{first-best (engagement on true class)} \\
  W_1 &\equiv W_0 - \text{exclusion externality}
       && \text{refused beneficial encounters} \\
  W_2 &\equiv W_1 - \text{coarseness loss}
       && \text{engagement on group standing, not class} \\
  W_3 &\equiv W_2 - \text{dynamic compounding}
       && \text{intergenerational drift in } \bar{c}_g \\
  W^{\mathrm{eq}} &\equiv W_3 - \text{platform cost}
       && \text{shadow market active.}
\end{align*}
Defining each loss component as a successive welfare gap,
\begin{equation}
  \mathcal{L}_{\text{alloc}} \equiv W_0 - W_1, \quad
  \mathcal{L}_{\text{info}} \equiv W_1 - W_2, \quad
  \mathcal{L}_{\text{dyn}} \equiv W_2 - W_3, \quad
  \mathcal{L}_{\text{platform}} \equiv W_3 - W^{\mathrm{eq}},
  \label{eq:welfare-defs}
\end{equation}
the total loss decomposes additively by telescoping:
\begin{equation}
  \mathcal{L} \;=\; W_0 - W^{\mathrm{eq}}
        \;=\; \mathcal{L}_{\text{alloc}} + \mathcal{L}_{\text{info}}
        + \mathcal{L}_{\text{dyn}} + \mathcal{L}_{\text{platform}}.
  \label{eq:welfare}
\end{equation}
Additivity is exact by construction (CHECK\_87) and each
$\mathcal{L}_k \geq 0$ because each counterfactual relaxation
removes a welfare-reducing distortion (CHECK\_88). The four
components map to model primitives as follows.

\subsection*{Allocative loss}
$\mathcal{L}_{\text{alloc}}$ is the expected foregone surplus from
refused beneficial encounters within a generation:
\begin{equation}
  \mathcal{L}_{\text{alloc}}
  \;=\; \int\!\!\!\int [\,U(c_i, c_j) - \bar{u}\,]_+
  \cdot \mathbf{1}\!\left[\,\E_j[d(c_i, c_j) \mid m_i] > \Delta^*(c_j)\,\right]
  \, dF(c_i)\, dF(c_j).
  \label{eq:L-alloc}
\end{equation}
Under the two-group engagement threshold of the main paper, this
reduces to $\mathcal{L}_{\text{alloc}} = s_g \cdot F(\bar{c}_g) \cdot
\overline{\Delta U}$, where $\overline{\Delta U}$ is the average
foregone surplus per blocked encounter and $s_g \cdot F(\bar{c}_g) = \phi$
is the population discrimination rate. The component is present
in \citet{Phelps1972}, \citet{Arrow1973}, and \citet{Loury2002}; it
is what standard models of statistical discrimination already
capture.

\subsection*{Informational loss}
$\mathcal{L}_{\text{info}}$ is the within-group foregone surplus that
arises because engagement conditions on the group standing
$\bar{c}_g$---all the marker conveys---rather than on the individual's
true class $c_i$. Writing $W^{\mathrm{class}}$ for the welfare under
true-class engagement and $W^{\mathrm{marker}}$ for the welfare under
group-standing engagement, both holding the engagement set fixed,
\begin{equation}
  \mathcal{L}_{\text{info}}
  \;=\; W^{\mathrm{class}} - W^{\mathrm{marker}}.
  \label{eq:L-info}
\end{equation}
This is the load-bearing object of the framework. A second-order
expansion of the bilateral surplus about the group mean gives a lower
bound proportional to the within-group class variance,
\begin{equation}
  \mathcal{L}_{\text{info}}
  \;\geq\; \tfrac{1}{2}\,\kappa\,\sigma^2_g \;>\; 0
  \qquad (\sigma^2_g > 0),
  \label{eq:L-info-bound}
\end{equation}
where $\sigma^2_g = \bar{c}_g(1-\bar{c}_g)/(1+\nu)$ is the within-group
class variance (with $\nu$ the concentration of the within-group class
distribution) and $\kappa = -\,\partial^2 U/\partial c_i^2 > 0$
measures the curvature of bilateral surplus in the partner's class. The
bound is \emph{independent of whether the standing $\bar{c}_g$ is
accurate}: it is generated by the coarseness of group-level
information, not by error in it. The loss falls on the whole society,
since group-standing engagement coarsens every interaction, not only
those involving the stigmatised group. There is no appeal here to a
Bayesian counterfactual: $\mathcal{L}_{\text{info}}$ is measured
against the true-class first-best $W^*$, and the model's
Jeffrey-versus-Bayes contrast governs belief \emph{dynamics}, not this
static floor.

\subsection*{Dynamic loss}
$\mathcal{L}_{\text{dyn}}$ is the discounted sum of allocative losses
across future generations attributable to the closure dynamic in
$\bar{c}_g$:
\begin{equation}
  \mathcal{L}_{\text{dyn}}
  \;=\; \sum_{t \geq 1} \beta^t \!\left[\,
  \mathcal{L}_{\text{alloc}}^{(t)} - \mathcal{L}_{\text{alloc}}^{(\mathrm{int})}
  \,\right]
  \label{eq:L-dyn}
\end{equation}
where $\mathcal{L}_{\text{alloc}}^{(\mathrm{int})}$ is the allocative
loss at the integration steady state $\bar{c}_g = 1$ (zero by
$F(1) = 0$). The component grows with the distance below the
bifurcation threshold $\gamma^*$ and falls primarily on future
generations of the stigmatised group.

\subsection*{Platform loss}
$\mathcal{L}_{\text{platform}}$ is the aggregate cost of the shadow
information market that arises when legal instruments constrain
formal screening:
\begin{equation}
  \mathcal{L}_{\text{platform}} \;=\; \int \Omega_j \, dj
  \label{eq:L-platform}
\end{equation}
the integral of shadow information value over all clients. Agents
invest resources in informal information---word of mouth, coded
reviews, network signals---that would be unnecessary in a
non-discriminatory equilibrium.

\section{Three regimes of welfare logic}
\label{sec:welfare-cases}

The reducible part of the welfare loss depends on what the group
standing $\bar{c}_g$ is doing relative to underlying empirical reality;
the within-group informational floor $\mathcal{L}_{\text{info}}$ is
present regardless. Three regimes exhaust the model's belief states,
each producing a distinct welfare logic and a distinct policy
comparison.

In the first regime, beliefs are accurate: $F_g$ tracks the empirical
class distribution within each group. The rigidity condition is
non-binding because there are no counterexamples to suppress, and
the closure trap does not arise as a structural feature. The model
behaves as a standard model of social exchange with an exclusion
externality.

In the second regime, beliefs are inaccurate but \emph{self-confirming}.
The conditionals $F_g$ do not match the empirical distribution, but
the equilibrium engagement decisions generate observed behaviour
that fails to falsify them: refused interactions produce no class
evidence, and accepted interactions selectively confirm the prior.
The rigidity condition prevents what little counter-evidence reaches
the agent from updating the conditional. Within a generation the
system is closed: beliefs produce behaviour that produces evidence
that fails to revise the beliefs.

In the third regime, beliefs are accurate and \emph{self-reinforcing}.
The system is at the closure equilibrium and the conditionals $F_g$
correctly describe the within-group distribution \emph{at} that
equilibrium. Beliefs match reality, but the equilibrium itself is
the trap: a higher-welfare equilibrium exists at $\bar{c}_g^H = 1$,
and the closure basin prevents the system from reaching it.
``Self-reinforcing'' here denotes intergenerational stability of the
equilibrium rather than within-generation stability of belief content;
the \emph{reducible} loss is in equilibrium selection, not in belief
error, while the within-group coarseness floor remains in force even at
the accurate standing.

These three regimes are exhaustive over the model's equilibrium
belief states: either beliefs track reality or they do not; if they
do not, the system either stabilises them through behaviour
(regime~2) or does not (a non-equilibrium configuration outside the
model's scope). Each regime corresponds to a different selection of
which \emph{reducible} components of $\mathcal{L}$ are non-zero---the
informational floor $\mathcal{L}_{\text{info}}$ is non-zero in all
three---and to a different policy logic.

\textbf{Case 1: Beliefs accurate.}
The standing $\bar{c}_g$ is correct, so there is no belief error to
compound: $\mathcal{L}_{\text{dyn}}$ vanishes once the system settles,
and the platform's shadow channel offers little advantage when beliefs
are already correct, so $\mathcal{L}_{\text{platform}}$ is small. Two
components remain. $\mathcal{L}_{\text{alloc}}$ is positive from the
exclusion externality that even accurate beliefs do not internalise,
and $\mathcal{L}_{\text{info}}$ is positive from the within-group
coarseness floor~(\ref{eq:L-info-bound})---engagement still runs on the
group standing, not the individual's class, so surplus is foregone even
though the standing is right (CHECK\_89). Policy: redistribution or
subsidised access compensates $\mathcal{L}_{\text{alloc}}$, but the
coarseness floor is untouched by any instrument that leaves markers
soft.

\textbf{Case 2: Beliefs inaccurate but self-confirming.}
All four components of $\mathcal{L}$ are positive (ANALYTIC\_11), and
the belief error in $\bar{c}_g$ inflates $\mathcal{L}_{\text{alloc}}$
and $\mathcal{L}_{\text{dyn}}$ above their accurate-standing levels.
This is the regime in which every instrument has a distinct incidence
across the components and the components interact most strongly. The
informational floor $\mathcal{L}_{\text{info}}$ is present here as in
every regime; what is special to Case~2 is the \emph{additional},
reducible loss that the inaccurate standing layers on top of it.

$\mathcal{L}_{\text{alloc}}$ responds to anti-discrimination law on
the formal-contracting side and to forced integration.
$\mathcal{L}_{\text{dyn}}$ responds to mobility programmes that
raise $\gamma$ toward the bifurcation threshold.
$\mathcal{L}_{\text{platform}}$ responds to platform regulation.
The fourth component, $\mathcal{L}_{\text{info}}$, is governed by
the irreducibility result (Corollary below).

\begin{corollary}[Irreducibility of informational welfare loss]
\label{cor:irreducibility-addendum}
Let $\Phi$ denote the class of policies that leave the soft-signal
structure of social interaction invariant: policies that operate by
changing $\gamma$ (mobility), $\sigma$ (admission), or $\delta$
(exclusion intensity), but not by converting soft markers into hard
signals. Then for all $\phi \in \Phi$, holding the group's internal
class dispersion fixed,
\[
  \mathcal{L}_{\mathrm{info}}(\phi)
  \;\geq\; \mathcal{L}_{\mathrm{info}}^*
  \;=\; \tfrac{1}{2}\,\kappa\,\sigma^2_g \;>\; 0
\]
whenever $\sigma^2_g > 0$. The floor is independent of whether the
standing $\bar{c}_g$ is accurate.
\end{corollary}

\begin{proof}[Proof sketch]
From~(\ref{eq:L-info-bound}), $\mathcal{L}_{\text{info}}$ is bounded
below by $\tfrac{1}{2}\kappa\,\sigma^2_g$, the second-order surplus
foregone when engagement conditions on the group mean rather than the
class realisation. Policies in $\Phi$ alter $\gamma$, $\sigma$, or
$\delta$; none of these enters the curvature $\kappa$ of bilateral
surplus or resolves the within-group class variance $\sigma^2_g$,
because each leaves the marker soft---informative about the group, not
the individual. The lower bound is therefore invariant to all
$\phi \in \Phi$ (CHECK\_86), and $\mathcal{L}_{\text{info}} > 0$
whenever the group carries internal class dispersion, $\sigma^2_g > 0$
(ANALYTIC\_10)---in particular at the accurate closure standing of
Case~3.
\qed
\end{proof}

This result characterises the limits of formal instruments within
the class $\Phi$. It does not preclude that social movements, norm
change, or long-run cultural shifts alter the signal structure
itself---moving markers from soft to hard, or reducing the
correlation between markers and group priors. Such changes would
operate outside $\Phi$ and could in principle reduce
$\mathcal{L}_{\text{info}}$. The corollary is a theorem about
instruments, not a verdict on the world.

\begin{table}[t]
\centering
\renewcommand{\arraystretch}{1.4}
\small
\begin{tabular}{lcccc}
\toprule
\textbf{Instrument} &
$\mathcal{L}_{\text{alloc}}$ &
$\mathcal{L}_{\text{info}}$ &
$\mathcal{L}_{\text{dyn}}$ &
$\mathcal{L}_{\text{platform}}$ \\
\midrule
Anti-discrimination law (formal contracting)
                          & $\downarrow$ & --- & --- & $\uparrow$ \\
Forced integration        & $\downarrow$ (temp.) & --- & --- & --- \\
Mobility programme ($\gamma > \gamma^*$) & $\downarrow\downarrow$ & --- & $\downarrow\downarrow$ & --- \\
Platform regulation       & --- & --- & --- & $\downarrow$ \\
\midrule
Irreducible               & & $\bullet$ & & \\
\bottomrule
\end{tabular}
\caption{Policy instruments mapped against welfare loss components
under Case~2. $\downarrow$ = reduces; $\uparrow$ = amplifies;
--- = no effect; $\bullet$ = structurally irreducible.
Anti-discrimination law on the formal-contracting side reduces
$\mathcal{L}_{\text{alloc}}$ on the contracting channel but
displaces activity into the informal contact channel, raising
$\mathcal{L}_{\text{platform}}$. Only the mobility programme
simultaneously eliminates allocative and dynamic loss above
$\gamma^*$.}
\label{tab:welfare-addendum}
\end{table}

\textbf{Case 3: Beliefs accurate and self-reinforcing.}
The standing $\bar{c}_g$ describes the realised closure equilibrium
correctly, yet $\mathcal{L}_{\text{info}} > 0$: even with an accurate
standing, engagement conditions on the group rather than the
individual, and the within-group coarseness
floor~(\ref{eq:L-info-bound}) is strictly positive whenever
$\sigma^2_g > 0$ (CHECK\_90). $\mathcal{L}_{\text{alloc}}$ remains
positive because the engagement threshold continues to refuse
beneficial encounters; $\mathcal{L}_{\text{dyn}}$ remains positive
because the trap reproduces itself across generations; and
$\mathcal{L}_{\text{platform}}$ remains positive because the shadow
market does not require belief error to function. This is the regime
the paper is built on---a homogeneous society at an accurate closure
standing---and it is precisely here that the irreducibility result has
force: the loss is structural, not a matter of belief content, and no
instrument in $\Phi$ removes it.

The relevant policy comparison is therefore across equilibria.
Forced integration achieves equilibrium selection by shifting a
cohort above $\bar{c}_g^{\mathrm{mid}}$, but does not alter the
trap's structure: subsequent cohorts face the same basin of
attraction. Mobility programmes raising $\gamma > \gamma^*$
eliminate the trap permanently through the saddle-node bifurcation;
above $\gamma^*$, the group converges from any starting position.
The two instruments are complements in the sub-threshold regime
($\gamma < \gamma^*$): integration manages cohort positions while
mobility investment pushes $\gamma$ toward $\gamma^*$. Integration
without mobility leaves the trap intact for future cohorts;
mobility above $\gamma^*$ is sufficient on its own. Platform
regulation reduces $\mathcal{L}_{\text{platform}}$ but has no
effect on the trap itself.

\section{Derivation of the welfare-dominance proposition}
\label{sec:welfare-dominance}

The main paper presents the following proposition as the headline
result of the welfare analysis:

\begin{proposition}[Welfare dominance of mobility above $\gamma^*$]
\label{prop:dominance}
Within the policy class $\Phi$, mobility programmes raising
$\gamma$ above $\gamma^*$ are uniquely welfare-dominant in the long
run: they eliminate $\mathcal{L}_{\mathrm{dyn}}$ permanently
through the saddle-node bifurcation, eliminate
$\mathcal{L}_{\mathrm{alloc}}$ at the integration steady state, and
do not amplify $\mathcal{L}_{\mathrm{platform}}$. No other
instrument in $\Phi$ achieves elimination on more than one
component. Anti-discrimination law on the formal-contracting side
addresses $\mathcal{L}_{\mathrm{alloc}}$ on that channel but
displaces activity to the informal contact channel, raising
$\mathcal{L}_{\mathrm{platform}}$; the net welfare effect is
ambiguous without empirical calibration.
\end{proposition}

\begin{proof}[Proof sketch]
Above $\gamma^*$ the saddle-node bifurcation destroys the closure
basin: the unique stable fixed point is the integration steady
state $\bar{c}_g = 1$. At this steady state, $F(\bar{c}_g) = 0$,
so the discrimination rate $\phi$ is zero and
$\mathcal{L}_{\text{alloc}}$ vanishes. With no further drift in
$\bar{c}_g$, $\mathcal{L}_{\text{dyn}}$ also vanishes. The
mobility instrument operates on $\gamma$, which does not enter the
platform's shadow-information optimisation, so
$\mathcal{L}_{\text{platform}}$ is unchanged. For any other
instrument $\phi \in \Phi$:
anti-discrimination law on the contracting side reduces
$\mathcal{L}_{\text{alloc}}$ on that channel but raises
$\mathcal{L}_{\text{platform}}$ via the shadow-market effect of
the main paper's Proposition~3.3.1; forced integration shifts a
cohort but does not eliminate the basin, so
$\mathcal{L}_{\text{dyn}}$ persists for subsequent cohorts;
platform regulation reduces $\mathcal{L}_{\text{platform}}$ only.
The intersection of components eliminated is therefore
$\{\mathcal{L}_{\text{alloc}}, \mathcal{L}_{\text{dyn}}\}$ for
mobility, a strict superset of what any alternative achieves.
\qed
\end{proof}

The dominance result, together with the irreducibility corollary,
produces a contract/contact reading of the policy implications.
Anti-discrimination law on the contract side
remains necessary on the formal contracting channel; mobility
investment is necessary on the contact channel; neither is
sufficient on its own, and at any interior standing both leave the
within-group informational floor $\mathcal{L}_{\text{info}}^*$
untouched, since neither converts the soft marker into a signal of
individual class. The result aligns with
\citet{Loury2002}: legal instruments address the contract channel,
not the contact channel; the contact channel runs through soft
signals that the model formalises via Jeffrey conditioning, and
its remedy requires instruments that operate on the
intergenerational class distribution rather than on individual
acts of exclusion.

\bibliographystyle{plainnat}
\bibliography{bibliography}

\end{document}
